\section{Interface}
\label{sec:interface}

\subsection{HTTP API}

The backend is served by FastAPI (\cref{tab:api}). Answering and retrieval
take the same request, a query and an optional document to scope it to. The
number of passages is deliberately not a parameter: it is decided by the
retrieval (\cref{sec:retrieval}).

\begin{table}[ht]
\centering
\small
\begin{tabularx}{\textwidth}{@{}llX@{}}
\toprule
Method & Path & Purpose \\
\midrule
GET    & \texttt{/health}                     & Liveness check. \\
GET    & \texttt{/stats}                      & CPU, GPU and memory load and
                                                temperatures. \\
GET    & \texttt{/documents}                  & Indexed documents and their
                                                chunk counts. \\
POST   & \texttt{/upload}                     & Store and ingest a file. \\
DELETE & \texttt{/document/\{name\}}          & Remove a document's chunks and
                                                cross-reference edges. \\
GET    & \texttt{/document/\{name\}/stats}    & Pages and structural units
                                                found. \\
POST   & \texttt{/query}                      & Retrieval-augmented answer. \\
POST   & \texttt{/retrieve}                   & Passages only, without
                                                generation. \\
GET    & \texttt{/document/\{name\}/provision}  & Full text of one provision,
                                                in reading order. \\
GET    & \texttt{/document/\{name\}/neighbours} & What a provision cites and
                                                what cites it. \\
\bottomrule
\end{tabularx}
\caption{API endpoints.}
\label{tab:api}
\end{table}

\subsection{Desktop application}

The interface is a React~19 application written in TypeScript and built with
Vite, running inside a Tauri~2 desktop shell. It is divided into three
resizable panes that can each be folded away (\cref{fig:interface}).

\begin{figure}[ht]
\centering
\IfFileExists{images/interface.png}{%
  \includegraphics[width=\textwidth]{images/interface.png}%
}{%
  \fbox{\parbox{0.9\textwidth}{\centering\small
  Screenshot of the application: save it as \texttt{images/interface.png}.}}%
}
\caption{The desktop application: retrieved citations (left), chat (centre)
and context (right).}
\label{fig:interface}
\end{figure}

\paragraph{Retrieved citations.} Each passage is shown as a card headed by its
source number, its printed page and its reference written the way it would be
cited---for example ``p.~118 $\cdot$ Article 101(1), second subparagraph''. A
line above the cards states how the passages were found: an exact match on the
provision named, or the closest passages by keyword and meaning. Below each
card, the provisions the passage cites can be followed (\cref{sec:graph}).

\paragraph{Chat.} The transcript shows the question, the answer, and the
detail beneath it in smaller type. Every answered turn keeps its own passages,
so clicking an earlier answer restores its citations without querying again. A
running query can be stopped. The composer shows which document questions are
scoped to and accepts a small set of commands, suggested as the first word is
typed: \texttt{find:} retrieves passages without generating an answer,
\texttt{doc:} changes the scope, \texttt{clear} empties the conversation and
\texttt{help} lists the commands. Shorthand for structural words is expanded
with the Tab key: \texttt{a}, \texttt{para}, \texttt{p}, \texttt{sec} and
\texttt{cha} become \emph{article}, \emph{paragraph}, \emph{point},
\emph{section} and \emph{chapter}. Only Tab expands, so the word ``a'' in an
ordinary sentence is left alone.

\paragraph{Context.} The last query's token usage is shown against the
4\,096-token window, split into prompt and response, together with the
session total and the time taken. Machine telemetry shows GPU and CPU load,
temperature and memory, and how much of the model is resident in video
memory. The document list shows each indexed document with a structural
summary (``144 pages $\cdot$ 13 chapters $\cdot$ 113 articles $\cdot$ 180
recitals''), accepts uploads of PDF, DOCX and TXT files with an estimate of
the indexing time, and scopes questions to a document when it is clicked.

\paragraph{Trade-offs.} Tauri was chosen over Electron because it uses the
operating system's web view instead of bundling a browser, which keeps the
application small and its memory footprint low beside a language model
competing for the same machine. The cost is that three web engines render the
interface---WebKit on macOS and Linux, WebView2 on Windows---and small
differences between them must be accommodated. Native translucency is used
only where the platform composites it itself (vibrancy on macOS, Mica and
Acrylic on Windows); elsewhere the window stays opaque, because a blur
redrawn by the web view on every change of the panes made the interface
sluggish.
